% Source: source/普通高考/2018/2018全国2文(甘肃,青海,内蒙古,黑龙江,吉林,辽宁,海南,宁夏,新疆,陕西,重庆).pdf, page 1.
% Type: flowchart (schematic; topology and labels are the mathematical content).
% Elements: 开始/结束 rounded boxes; N=0,T=0, i=1, N=N+1/i,
% T=T+1/(i+1), blank increment box, and S=N-T process boxes; decision i<100;
% 输出 S parallelogram; directed connectors and 是/否 branch labels.
% Relations: 开始 -> N,T initialization -> i=1 -> i<100; 是 loops through
% the two accumulators and the blank increment box back to the decision;
% 否 -> S=N-T -> 输出 S -> 结束.  All connectors and box borders are solid.
% Node-edge checklist transcribed from the PDF: 开始→N=0,T=0; N=0,T=0→i=1;
% i=1→i<100; i<100(是)→N=N+1/i; N=N+1/i→T=T+1/(i+1);
% T=T+1/(i+1)→空白框; 空白框→i<100; i<100(否)→S=N−T;
% S=N−T→输出S; 输出S→结束.  Each listed edge is drawn exactly once below.
% Label strategy: branch labels sit outside the diamond; math labels are centered
% in their boxes; the return path stays left of the processing column.

\begin{tikzpicture}[
  >=Stealth,
  line width=0.45pt,
  node distance=0pt,
  flowtext/.style={anchor=center, align=center,
    execute at begin node={%
      \everypar{}%
      \setlength{\parindent}{0pt}%
      \setlength{\leftskip}{0pt}%
      \setlength{\rightskip}{0pt}%
    }},
  every node/.style={flowtext, font=\normalsize},
  flowbox/.style={draw, flowtext, minimum height=0.72cm,
    inner xsep=4pt, inner ysep=2.5pt},
  rounded/.style={flowbox, rounded corners=3pt},
  decision/.style={draw, flowtext, diamond, aspect=2.7, minimum width=3.5cm,
    minimum height=1.0cm},
  io/.style={draw, flowtext, trapezium, trapezium left angle=74,
    trapezium right angle=106, minimum height=0.72cm, minimum width=2.3cm}
]
  % Main vertical sequence.
  \node[rounded, minimum width=1.15cm] (start) at (0,11.7) {开始};
  \node[flowbox, minimum width=3.8cm] (init) at (0,10.45) {$N=0,T=0$};
  \node[flowbox, minimum width=1.7cm] (seti) at (0,9.35) {$i=1$};
  \node[decision] (test) at (0,8.08) {$i<100$};

  \node[flowbox, minimum width=2.55cm] (accn) at (-2.0,6.65)
    {$N=N+\dfrac{1}{i}$};
  \node[flowbox, minimum width=3.0cm] (acct) at (-2.0,5.18)
    {$T=T+\dfrac{1}{i+1}$};
  \node[flowbox, minimum width=2.05cm, minimum height=0.72cm] (inc) at (-2.0,3.65) {};

  \node[flowbox, minimum width=2.75cm] (sub) at (2.0,6.65) {$S=N-T$};
  \node[io, minimum width=2.4cm] (out) at (2.0,5.15) {输出 $S$};
  \node[rounded, minimum width=1.15cm] (finish) at (2.0,3.75) {结束};

  % Main arrows.
  \draw[->] (start.south) -- (init.north);
  \draw[->] (init.south) -- (seti.north);
  % The source has a plain trunk down to the loop re-entry, then a separate
  % downward arrow into the decision diamond (no overlapping arrowheads).
  \coordinate (trunkjoin) at (0,8.40);
  \coordinate (testentry) at (0.80,8.40);
  \draw (seti.south) -- (trunkjoin);
  \draw (trunkjoin) -- (testentry);
  \draw[->] (testentry) -- (test.north);
  \draw[->] (accn.south) -- (acct.north);
  \draw[->] (acct.south) -- (inc.north);
  \draw[->] (test.east) -- node[above] {否} (2.0,8.08) -- (sub.north);
  \draw[->] (sub.south) -- (out.north);
  \draw[->] (out.south) -- (finish.north);

  % Yes branch and loop-back path.
  \draw[->] (test.west) -- node[above] {是} (-2.0,8.08) -- (accn.north);
  % The return arrow first joins the main trunk with a right-pointing head;
  % the separate trunk arrow above then enters the decision diamond.
  \draw[->] (inc.west) -- (-3.90,3.65) -- (-3.90,8.40) -- (-1.00,8.40);
  \draw (-1.00,8.40) -- (trunkjoin);
\end{tikzpicture}
